\scanpage{73}
% archive foliation: 72
et $S,\mathcal J$ sont deux parties discrètes disjointes de $\widehat X$.

On va \deleted{\unsure{se}} passer aux catégories isotopiques correspondantes, et en donner des interprétations en termes de théorie des groupes. Bien entendu, pour cette description, on est ramené au cas connexe. On peut \deleted{aussi} noter tt de suite que pour ce pb, on est ramené au cas où $\mathcal J=\varnothing$ [en \deleted{remplaçant}\inserted{\unsure{décrivant|donnant}} dans le cas général comme celui de la donnée d'un $(\widehat X,I)$ \uline{plus} la donnée d'une partie $\mathcal J$ de $I$]. Mais on gardera le pb \unsure{ss|sous} \inserted{la} forme générale, en gardant également la dissymétrie dans le rôle de $S,\mathcal J$ -- mais
